\documentclass[a4paper,11pt,oneside]{memoir}

% Castellano
\usepackage[spanish,es-tabla]{babel}
\selectlanguage{spanish}
\usepackage[utf8]{inputenc}
\usepackage{placeins}

\RequirePackage{booktabs}
\RequirePackage[table]{xcolor}
\RequirePackage{xtab}
\RequirePackage{multirow}

% Links
\usepackage[colorlinks]{hyperref}
\hypersetup{
	allcolors = {red}
}

% Ecuaciones
\usepackage{amsmath}

% Rutas de fichero / paquete
\newcommand{\ruta}[1]{{\sffamily #1}}

% Párrafos
\nonzeroparskip


% Imagenes
\usepackage{graphicx}
\newcommand{\imagen}[2]{
	\begin{figure}[!h]
		\centering
		\includegraphics[width=0.9\textwidth]{#1}
		\caption{#2}\label{fig:#1}
	\end{figure}
	\FloatBarrier
}

\newcommand{\imagenflotante}[2]{
	\begin{figure}%[!h]
		\centering
		\includegraphics[width=0.9\textwidth]{#1}
		\caption{#2}\label{fig:#1}
	\end{figure}
}



% El comando \figura nos permite insertar figuras comodamente, y utilizando
% siempre el mismo formato. Los parametros son:
% 1 -> Porcentaje del ancho de página que ocupará la figura (de 0 a 1)
% 2 --> Fichero de la imagen
% 3 --> Texto a pie de imagen
% 4 --> Etiqueta (label) para referencias
% 5 --> Opciones que queramos pasarle al \includegraphics
% 6 --> Opciones de posicionamiento a pasarle a \begin{figure}
\newcommand{\figuraConPosicion}[6]{%
  \setlength{\anchoFloat}{#1\textwidth}%
  \addtolength{\anchoFloat}{-4\fboxsep}%
  \setlength{\anchoFigura}{\anchoFloat}%
  \begin{figure}[#6]
    \begin{center}%
      \Ovalbox{%
        \begin{minipage}{\anchoFloat}%
          \begin{center}%
            \includegraphics[width=\anchoFigura,#5]{#2}%
            \caption{#3}%
            \label{#4}%
          \end{center}%
        \end{minipage}
      }%
    \end{center}%
  \end{figure}%
}

%
% Comando para incluir imágenes en formato apaisado (sin marco).
\newcommand{\figuraApaisadaSinMarco}[5]{%
  \begin{figure}%
    \begin{center}%
    \includegraphics[angle=90,height=#1\textheight,#5]{#2}%
    \caption{#3}%
    \label{#4}%
    \end{center}%
  \end{figure}%
}
% Para las tablas
\newcommand{\otoprule}{\midrule [\heavyrulewidth]}
%
% Nuevo comando para tablas pequeñas (menos de una página).
\newcommand{\tablaSmall}[5]{%
 \begin{table}
  \begin{center}
   \rowcolors {2}{gray!35}{}
   \begin{tabular}{#2}
    \toprule
    #4
    \otoprule
    #5
    \bottomrule
   \end{tabular}
   \caption{#1}
   \label{tabla:#3}
  \end{center}
 \end{table}
}

%
% Nuevo comando para tablas pequeñas (menos de una página).
\newcommand{\tablaSmallSinColores}[5]{%
 \begin{table}[H]
  \begin{center}
   \begin{tabular}{#2}
    \toprule
    #4
    \otoprule
    #5
    \bottomrule
   \end{tabular}
   \caption{#1}
   \label{tabla:#3}
  \end{center}
 \end{table}
}

\newcommand{\tablaApaisadaSmall}[5]{%
\begin{landscape}
  \begin{table}
   \begin{center}
    \rowcolors {2}{gray!35}{}
    \begin{tabular}{#2}
     \toprule
     #4
     \otoprule
     #5
     \bottomrule
    \end{tabular}
    \caption{#1}
    \label{tabla:#3}
   \end{center}
  \end{table}
\end{landscape}
}

%
% Nuevo comando para tablas grandes con cabecera y filas alternas coloreadas en gris.
\newcommand{\tabla}[6]{%
  \begin{center}
    \tablefirsthead{
      \toprule
      #5
      \otoprule
    }
    \tablehead{
      \multicolumn{#3}{l}{\small\sl continúa desde la página anterior}\\
      \toprule
      #5
      \otoprule
    }
    \tabletail{
      \hline
      \multicolumn{#3}{r}{\small\sl continúa en la página siguiente}\\
    }
    \tablelasttail{
      \hline
    }
    \bottomcaption{#1}
    \rowcolors {2}{gray!35}{}
    \begin{xtabular}{#2}
      #6
      \bottomrule
    \end{xtabular}
    \label{tabla:#4}
  \end{center}
}

%
% Nuevo comando para tablas grandes con cabecera.
\newcommand{\tablaSinColores}[6]{%
  \begin{center}
    \tablefirsthead{
      \toprule
      #5
      \otoprule
    }
    \tablehead{
      \multicolumn{#3}{l}{\small\sl continúa desde la página anterior}\\
      \toprule
      #5
      \otoprule
    }
    \tabletail{
      \hline
      \multicolumn{#3}{r}{\small\sl continúa en la página siguiente}\\
    }
    \tablelasttail{
      \hline
    }
    \bottomcaption{#1}
    \begin{xtabular}{#2}
      #6
      \bottomrule
    \end{xtabular}
    \label{tabla:#4}
  \end{center}
}

%
% Nuevo comando para tablas grandes sin cabecera.
\newcommand{\tablaSinCabecera}[5]{%
  \begin{center}
    \tablefirsthead{
      \toprule
    }
    \tablehead{
      \multicolumn{#3}{l}{\small\sl continúa desde la página anterior}\\
      \hline
    }
    \tabletail{
      \hline
      \multicolumn{#3}{r}{\small\sl continúa en la página siguiente}\\
    }
    \tablelasttail{
      \hline
    }
    \bottomcaption{#1}
  \begin{xtabular}{#2}
    #5
   \bottomrule
  \end{xtabular}
  \label{tabla:#4}
  \end{center}
}



\definecolor{cgoLight}{HTML}{EEEEEE}
\definecolor{cgoExtralight}{HTML}{FFFFFF}

%
% Nuevo comando para tablas grandes sin cabecera.
\newcommand{\tablaSinCabeceraConBandas}[5]{%
  \begin{center}
    \tablefirsthead{
      \toprule
    }
    \tablehead{
      \multicolumn{#3}{l}{\small\sl continúa desde la página anterior}\\
      \hline
    }
    \tabletail{
      \hline
      \multicolumn{#3}{r}{\small\sl continúa en la página siguiente}\\
    }
    \tablelasttail{
      \hline
    }
    \bottomcaption{#1}
    \rowcolors[]{1}{cgoExtralight}{cgoLight}

  \begin{xtabular}{#2}
    #5
   \bottomrule
  \end{xtabular}
  \label{tabla:#4}
  \end{center}
}




\graphicspath{ {./img/} }

% Capítulos
\chapterstyle{bianchi}
\newcommand{\capitulo}[2]{
	\setcounter{chapter}{#1}
	\setcounter{section}{0}
	\chapter*{#2}
	\addcontentsline{toc}{chapter}{#2}
	\markboth{#2}{#2}
}

% Apéndices
\renewcommand{\appendixname}{Apéndice}
\renewcommand*\cftappendixname{\appendixname}

\newcommand{\apendice}[1]{
	%\renewcommand{\thechapter}{A}
	\chapter{#1}
}

\renewcommand*\cftappendixname{\appendixname\ }

% Formato de portada
\makeatletter
\usepackage{xcolor}
\newcommand{\tutor}[1]{\def\@tutor{#1}}
\newcommand{\course}[1]{\def\@course{#1}}
\definecolor{cpardoBox}{HTML}{E6E6FF}
\def\maketitle{
  \null
  \thispagestyle{empty}
  % Cabecera ----------------
\noindent\includegraphics[width=\textwidth]{cabecera}\vspace{1cm}%
  \vfill
  % Título proyecto y escudo informática ----------------
  \colorbox{cpardoBox}{%
    \begin{minipage}{.8\textwidth}
      \vspace{.5cm}\Large
      \begin{center}
      \textbf{TFG del Grado en Ingeniería Informática}\vspace{.6cm}\\
      \textbf{\LARGE\@title{}}
      \end{center}
      \vspace{.2cm}
    \end{minipage}

  }%
  \hfill\begin{minipage}{.20\textwidth}
    \includegraphics[width=\textwidth]{escudoInfor}
  \end{minipage}
  \vfill
  % Datos de alumno, curso y tutores ------------------
  \begin{center}%
  {%
    \noindent\LARGE
    Presentado por \@author{}\\ 
    en Universidad de Burgos --- \@date{}\\
    Tutor: \@tutor{}\\
  }%
  \end{center}%
  \null
  \cleardoublepage
  }
\makeatother


% Datos de portada
\title{título del TFG \\Documentación Técnica}
\author{nombre alumno}
\tutor{nombre tutor}
\date{\today}

\begin{document}

\maketitle



\cleardoublepage



%%%%%%%%%%%%%%%%%%%%%%%%%%%%%%%%%%%%%%%%%%%%%%%%%%%%%%%%%%%%%%%%%%%%%%%%%%%%%%%%%%%%%%%%



\frontmatter


\clearpage

% Indices
\tableofcontents

\clearpage

\listoffigures

\clearpage

\listoftables

\clearpage

\mainmatter

\appendix

\apendice{Plan de Proyecto Software}

\section{Introducción}

En este anexo se tratan dos puntos principales:

\begin{itemize}
	\item Planificación temporal del proyecto, donde se han utilizado metodologías ágiles.
	\item Estudio de la viabilidad del proyecto tanto a nivel económico como legal.
\end{itemize}

\section{Planificación temporal}

La planificación temporal del proyecto se caracteriza por el uso de metodologías ágiles, concretamente Scrum. El desarrollo se ha dividido en una serie de iteraciones (o \textit{sprints}).

Se han realizado un total de 14 \textit{sprints}. Cada uno de ellos está enfocado a una serie de objetivos particulares. Cada \textit{sprint} se corresponde con una \textit{milestone} del repositorio en GitHub. A su vez, cada \textit{milestone} está compuesta por una serie de tareas (o \textit{issues}).

Cada una de las iteraciones finaliza con una reunión presencial o mediante videoconferencia donde se revisan y aceptan los cambios. Tras ello, los cambios aceptados son integrados en la rama principal (rama \textit{master}).

El \textit{add-on} ZenHub se ha utilizado para asignar estimaciones temporales a las tareas, las cuales permiten el desarrollo de un diagrama \textit{burndown} por \textit{sprint}. Para las estimaciones de tiempo se utiliza la unidad \textit{story point}, a continuación se muestra una tabla de equivalencias aproximadas:

\tablaSmallSinColores{Equivalencia \textit{story-point} a tiempo.}{l l}{story-point}
{ \textit{Story points} & Tiempo \\}{ 
1 & 30 minutos \\
2 & 1 hora \\
3 & 2 horas \\
4 & 3 horas \\
5 & 6 horas \\
6 & 8 horas \\
8 & 10 horas \\
13 & 16 horas \\
}

A continuación se muestran las iteraciones del proyecto:

\subsection{\textit{Sprint} 1 (10/02/2017 - 24/02/2017)}

Esta iteración del proyecto contiene tareas relacionadas con el aprendizaje de conceptos y herramientas. Sus dos tareas suman una estimación de 20 \textit{story points}.

\url{https://github.com/mjuez/TFM2016_Analisis-Visual-Revisiones-Codigo/milestone/1}

\imagen{sprint1}{Diagrama \textit{burndown} del \textit{sprint} 1.}

\subsection{\textit{Sprint} 2 (24/02/2017 - 10/03/2017)}

Esta iteración del proyecto contiene tareas relacionadas con el desarrollo de la memoria y prototipos previos. Sus 4 tareas suman una estimación de 22 \textit{story points}.

\url{https://github.com/mjuez/TFM2016_Analisis-Visual-Revisiones-Codigo/milestone/2}

\imagen{sprint2}{Diagrama \textit{burndown} del \textit{sprint} 2.}

\subsection{\textit{Sprint} 3 (10/03/2017 - 24/03/2017)}

Esta iteración del proyecto contiene tareas relacionadas con el desarrollo de la memoria, prototipos previos y estudio de tecnologías a utilizar. Sus 5 tareas suman una estimación de 17 \textit{story points}.

\url{https://github.com/mjuez/TFM2016_Analisis-Visual-Revisiones-Codigo/milestone/3}

\imagen{sprint3}{Diagrama \textit{burndown} del \textit{sprint} 3.}

\subsection{\textit{Sprint} 4 (24/03/2017 - 07/04/2017)}

Esta iteración del proyecto contiene tareas relacionadas con el desarrollo de la memoria, prototipos previos y estudio de herramientas de revisión automática. Sus 6 tareas suman una estimación de 18 \textit{story points}.

\url{https://github.com/mjuez/TFM2016_Analisis-Visual-Revisiones-Codigo/milestone/4}

\imagen{sprint4}{Diagrama \textit{burndown} del \textit{sprint} 4.}

\subsection{\textit{Sprint} 5 (07/04/2017 - 21/04/2017)}

Esta iteración del proyecto contiene tareas relacionadas con el estudio de la API de GitHub y el desarrollo de la parte \textit{backend} de la aplicación final. Sus 8 tareas suman una estimación de 32 \textit{story points}.

\url{https://github.com/mjuez/TFM2016_Analisis-Visual-Revisiones-Codigo/milestone/5}

\imagen{sprint5}{Diagrama \textit{burndown} del \textit{sprint} 5.}

\subsection{\textit{Sprint} 6 (21/04/2017 - 28/04/2017)}

Esta iteración del proyecto contiene tareas relacionadas con el estudio de la de las revisiones en GitHub, desarrollo de documentación, desarrollo de la parte \textit{backend} de la aplicación final, y creación de pruebas unitarias. Sus 7 tareas suman una estimación de 24 \textit{story points}.

\url{https://github.com/mjuez/TFM2016_Analisis-Visual-Revisiones-Codigo/milestone/6}

\imagen{sprint6}{Diagrama \textit{burndown} del \textit{sprint} 6.}

\subsection{\textit{Sprint} 7 (28/04/2017 - 12/05/2017)}

Esta iteración del proyecto contiene tareas relacionadas con el desarrollo y finalización de la parte \textit{backend} de la aplicación final. Sus 9 tareas suman una estimación de 34 \textit{story points}.

\url{https://github.com/mjuez/TFM2016_Analisis-Visual-Revisiones-Codigo/milestone/7}

\imagen{sprint7}{Diagrama \textit{burndown} del \textit{sprint} 7.}

\subsection{\textit{Sprint} 8 (12/05/2017 - 19/05/2017)}

Esta iteración del proyecto contiene tareas relacionadas con refactorizaciones de la parte \textit{backend} e inicio de la parte \textit{frontend}. Sus 4 tareas suman una estimación de 20 \textit{story points}.

\url{https://github.com/mjuez/TFM2016_Analisis-Visual-Revisiones-Codigo/milestone/8}

\imagen{sprint8}{Diagrama \textit{burndown} del \textit{sprint} 8.}

\subsection{\textit{Sprint} 9 (19/05/2017 - 26/05/2017)}

Esta iteración del proyecto contiene tareas relacionadas con el desarrollo de la parte \textit{frontend} y alguna mejora en la parte \textit{backend}. Sus 5 tareas suman una estimación de 20 \textit{story points}.

\url{https://github.com/mjuez/TFM2016_Analisis-Visual-Revisiones-Codigo/milestone/9}

\imagen{sprint9}{Diagrama \textit{burndown} del \textit{sprint} 9.}

\subsection{\textit{Sprint} 10 (26/05/2017 - 02/06/2017)}

Esta iteración del proyecto contiene tareas relacionadas con el desarrollo de la parte \textit{frontend}. Sus 6 tareas suman una estimación de 33 \textit{story points}.

\url{https://github.com/mjuez/TFM2016_Analisis-Visual-Revisiones-Codigo/milestone/10}

\imagen{sprint10}{Diagrama \textit{burndown} del \textit{sprint} 10.}

\subsection{\textit{Sprint} 11 (02/06/2017 - 09/06/2017)}

Esta iteración del proyecto contiene tareas relacionadas con la finalización del desarrollo de la parte \textit{frontend} y alguna mejora en la parte \textit{backend}. Sus 8 tareas suman una estimación de 23 \textit{story points}.

\url{https://github.com/mjuez/TFM2016_Analisis-Visual-Revisiones-Codigo/milestone/11}

\imagen{sprint11}{Diagrama \textit{burndown} del \textit{sprint} 11.}

\subsection{\textit{Sprint} 12 (09/06/2017 - 19/06/2017)}

Esta iteración del proyecto contiene tareas relacionadas con la implementación de las últimas características de la aplicación final, documentación de código, y desarrollo de pruebas unitarias. Sus 7 tareas suman una estimación de 34 \textit{story points}.

\url{https://github.com/mjuez/TFM2016_Analisis-Visual-Revisiones-Codigo/milestone/12}

\imagen{sprint12}{Diagrama \textit{burndown} del \textit{sprint} 12.}

\subsection{\textit{Sprint} 13 (19/06/2017 - 26/06/2017)}

Esta iteración del proyecto contiene tareas relacionadas con la finalización del desarrollo de la memoria del proyecto. Sus 11 tareas suman una estimación de 44 \textit{story points}.

\url{https://github.com/mjuez/TFM2016_Analisis-Visual-Revisiones-Codigo/milestone/13}

\imagen{sprint13}{Diagrama \textit{burndown} del \textit{sprint} 13.}

\subsection{\textit{Sprint} 14 (26/06/2017 - 03/07/2017)}

Ésta es la última iteración del proyecto. Sus tareas están relacionadas con el desarrollo de los anexos (documentación técnica).

\url{https://github.com/mjuez/TFM2016_Analisis-Visual-Revisiones-Codigo/milestone/14}

En el momento de realización de este documento, la iteración aún no había terminado, y por tanto no existía el diagrama \textit{burndown}.


\section{Estudio de viabilidad}

A continuación se desarrolla el estudio de viabilidad del proyecto, tanto a nivel económico como a nivel legal.

\subsection{Análisis coste-beneficio}

El estudio de viabilidad económica consiste en el análisis de costes y beneficios del proyecto.

\subsubsection{Costes}

Los costes del proyecto se han desglosado en los siguientes tipos: \newline

\textbf{Costes de personal:}

El proyecto ha sido realizado por un desarrollador empleado a media jornada durante 3 meses, y a jornada completa durante 2 meses:

\tablaSmallConPie{Costes de personal a media jornada.}{l l}{costes-personal-mj}
{ \textbf{Descripción} & \textbf{Coste} \\}{ 
Salario mensual neto & 550,00 \euro{} \\
Retención IRPF (19\%) & 204,50 \euro{} \\
Seguridad Social (29,9\%) & 321,82 \euro{} \\
Salario mensual bruto & 1.076,32 \euro{} \\
}{
\textbf{Total 3 meses} & 3.228,96 \euro{} \\
}

\tablaSmallConPie{Costes de personal a jornada completa.}{l l}{costes-personal-jc}
{ \textbf{Descripción} & \textbf{Coste} \\}{ 
Salario mensual neto & 1.100,00 \euro{} \\
Retención IRPF (19\%) & 409,00 \euro{} \\
Seguridad Social (29,9\%) & 643,64 \euro{} \\
Salario mensual bruto & 2.152,64 \euro{} \\
}{
\textbf{Total 2 meses} & 4.305,28 \euro{} \\
}

Los costes de Seguridad Social (29,9\%) se han calculado tal como indica el Ministerio de Empleo y Seguridad Social \cite{seguridadsocial}:
\begin{itemize}
	\item 23,60\% por contingencias comunes.
	\item 5,50\% por desempleo de tipo general.
	\item 0,20\% para el Fondo de Garantía Salarial.
	\item 0,60\% por formación profesional.
\end{itemize}

\textbf{Costes de \textit{hardware}:}

Para el desarrollo del proyecto se ha utilizado un ordenador portátil con un coste de 990 \euro{} y una vida útil de 5 años. El coste amortizado es de 82,5 \euro.

Para calcular el coste de amortización se ha utilizado la siguiente fórmula:

$$ Coste Amortizado = \frac{Coste Hardware}{Vida Util} * Tiempo Utilizado $$ \newline

\textbf{Costes de \textit{software}:}

Todo el \textit{software} utilizado durante el desarrollo del proyecto es gratuito o tiene planes gratuitos. \newline

\textbf{Otros costes:}

A continuación se detalla el resto de costes del proyecto:

\tablaSmallConPie{Otros costes.}{l l}{otros-costes}
{ \textbf{Descripción} & \textbf{Coste} \\}{ 
Internet & 165,25 \euro{} \\
Documentación (memoria impresa + \textit{CD's}) & 50,00 \euro{} \\
}{
\textbf{Total} & 215,25 \euro{} \\
}

\textbf{Costes totales:}

La suma de todos los costes del proyecto es la siguiente:

\tablaSmallConPie{Costes totales.}{l l}{costes-totales}
{ \textbf{Descripción} & \textbf{Coste} \\}{ 
Personal & 7.534,24 \euro{} \\
\textit{Hardware} & 82,50 \euro{} \\
\textit{Software} & 0,00 \euro{} \\
Otros & 215,25 \euro{} \\
}{
\textbf{Total} & 7.831,99 \euro{} \\
}

\subsubsection{Beneficios}

Este proyecto se ha desarrollado para ser utilizado en entornos de investigación universitaria. No se espera obtener beneficio económico del mismo.

\subsection{Viabilidad legal}

El estudio de la viabilidad legal se ha realizado en términos relativos a las licencias utilizadas.

\subsubsection{Licencia del \textit{software} desarrollado}

Como se ha indicado anteriormente, nuestro proyecto tiene fines académicos y de investigación. Por ello no queremos restringir su uso, distribución o modificación.

Para la elección de la licencia de la aplicación desarrollada se han tenido en cuenta las licencias de las dependencias externas utilizadas.

\tablaSmallSinColores{Diferentes licencias utilizadas por nuestras dependencias.}{l l}{licencias}
{ \textbf{Licencia} & \textbf{Tipo} \\}{ 
WTFPL & Copyright \\
MIT & Copyright \\
ISC & Copyright \\
Apache 2.0 & Copyright \\
BSD-3-Clause & Copyright \\
}

La licencia MIT \cite{mit} se ajusta a nuestras necesidades. Es compatible con todas las anteriores, y además es poco restrictiva. Permite, entre otras cosas, el uso, copia, modificación o distribución del \textit{software} con la única condición de incluir nuestro aviso de copyright y licencia en toda copia o derivado.

La licencia GPLv3 también era compatible con las licencias de nuestras dependencias, pero al ser de tipo copyleft, es algo más restrictiva que la MIT.

\subsubsection{Licencia de la documentación}

En lo relativo a la documentación (memoria y anexos), se ha utilizado una licencia de tipo Creative Commons 4.0 \cite{creativecommons} (CC BY 4.0), que permite la copia y redistribución, así como su modificación para cualquier propósito.
\apendice{Especificación de Requisitos}

\section{Introducción}

\section{Objetivos generales}

\section{Catalogo de requisitos}

\section{Especificación de requisitos}



\apendice{Especificación de diseño}

\section{Introducción}

\section{Diseño de datos}

\section{Diseño procedimental}

\section{Diseño arquitectónico}



\apendice{Documentación técnica de programación}

\section{Introducción}

\section{Estructura de directorios}

\section{Manual del programador}

\section{Compilación, instalación y ejecución del proyecto}

\section{Pruebas del sistema}

\apendice{Documentación de usuario}

\section{Introducción}

\section{Requisitos de usuarios}

\section{Instalación}

\section{Manual del usuario}





\bibliographystyle{plain}
\bibliography{bibliografiaAnexos}

\end{document}
